% chapters/06/latihan.tex
\bagian{Latihan dan Refleksi}

\begin{enumerate}
  \item Tuliskan dua baris perhitungan yang memperbarui kecepatan lalu posisi
        pada integrasi Euler, dan jelaskan mengapa urutannya tidak boleh
        dibalik.

  \item Jelaskan mengapa percepatan gravitasi bernilai positif pada koordinat
        layar, sedangkan kecepatan lompatan bernilai negatif.

  \item Sebuah benda bergerak dengan kecepatan 300 piksel per detik ke bawah
        dan gravitasi 1500 piksel per detik kuadrat pada laju 60 bingkai per
        detik. Hitunglah kecepatannya pada bingkai berikutnya dan perpindahan
        yang terjadi pada bingkai itu.

  \item Jelaskan mengapa pemain sebaiknya diberi koefisien restitusi nol,
        sedangkan bola tidak. Apa yang terjadi apabila keduanya ditukar?

  \item Sebuah lompatan hendak dibuat setinggi 200 piksel. Hitunglah kecepatan
        lompatan yang diperlukan apabila gravitasinya 1800 piksel per detik
        kuadrat, lalu hitunglah lama benda berada di udara.

  \item Jelaskan mengapa lompatan harus dibaca sebagai peristiwa
        \code{KEYDOWN}, bukan sebagai keadaan tombol dengan
        \code{key.get_pressed()}. Sebutkan gejala yang muncul apabila pilihan
        itu tertukar.

  \item Sebuah bola dijatuhkan dari ketinggian 100 piksel dengan koefisien
        restitusi 0,7. Hitunglah tinggi pantulan pertama dan kedua, lalu
        jelaskan mengapa tingginya menyusut secepat itu.

  \item Perhatikan potongan kode berikut, lalu tentukan kesalahannya dan
        perbaikannya.
\begin{kodekeluaran}
if position.y + half == floor_y:
    velocity.y = -velocity.y * restitution
    position.y = floor_y - half
\end{kodekeluaran}
\end{enumerate}

\begin{refleksibox}
\begin{enumerate}
  \item Pada langkah keempat praktikum, pasangan nilai mana yang menurut Anda
        paling enak dimainkan? Mengapa demikian, dan bagaimana jawaban Anda
        berhubungan dengan lama benda berada di udara?
  \item Ketika koefisien restitusi bola diubah menjadi 1,0, apa yang terjadi?
        Adakah alasannya untuk memakai nilai itu pada permainan yang
        sesungguhnya?
  \item Pada langkah keenam praktikum, pemain hanya dapat mendarat di atas
        bongkahan tetapi tidak dapat menabraknya dari samping. Menurut Anda,
        pemeriksaan seperti apa yang diperlukan untuk menanganinya?
\end{enumerate}
\end{refleksibox}
